\documentclass[a4paper,12pt,twoside]{memoir}

% Castellano
\usepackage[spanish,es-tabla]{babel}
\selectlanguage{spanish}
\usepackage[utf8]{inputenc}
\usepackage[T1]{fontenc}
\usepackage{lmodern} % Scalable font
\usepackage{microtype}
\usepackage{placeins}

\usepackage[backend=bibtex,dateabbrev=false]{biblatex}
\addbibresource{bibliografia.bib}

\RequirePackage{booktabs}
\RequirePackage[table]{xcolor}
\RequirePackage{xtab}
\RequirePackage{multirow}

% Links
\PassOptionsToPackage{hyphens}{url}\usepackage[colorlinks]{hyperref}
\hypersetup{
	allcolors = {red}
}

% Ecuaciones
\usepackage{amsmath}

% Rutas de fichero / paquete
\newcommand{\ruta}[1]{{\sffamily #1}}

% Párrafos
\nonzeroparskip

% Huérfanas y viudas
\widowpenalty100000
\clubpenalty100000

% Imagenes
\usepackage{graphicx}
\newcommand{\imagen}[2]{
	\begin{figure}[!h]
		\centering
		\includegraphics[width=0.9\textwidth]{#1}
		\caption{#2}\label{fig:#1}
	\end{figure}
	\FloatBarrier
}

\newcommand{\imagenflotante}[2]{
	\begin{figure}%[!h]
		\centering
		\includegraphics[width=0.9\textwidth]{#1}
		\caption{#2}\label{fig:#1}
	\end{figure}
}



% El comando \figura nos permite insertar figuras comodamente, y utilizando
% siempre el mismo formato. Los parametros son:
% 1 -> Porcentaje del ancho de página que ocupará la figura (de 0 a 1)
% 2 --> Fichero de la imagen
% 3 --> Texto a pie de imagen
% 4 --> Etiqueta (label) para referencias
% 5 --> Opciones que queramos pasarle al \includegraphics
% 6 --> Opciones de posicionamiento a pasarle a \begin{figure}
\newcommand{\figuraConPosicion}[6]{%
  \setlength{\anchoFloat}{#1\textwidth}%
  \addtolength{\anchoFloat}{-4\fboxsep}%
  \setlength{\anchoFigura}{\anchoFloat}%
  \begin{figure}[#6]
    \begin{center}%
      \Ovalbox{%
        \begin{minipage}{\anchoFloat}%
          \begin{center}%
            \includegraphics[width=\anchoFigura,#5]{#2}%
            \caption{#3}%
            \label{#4}%
          \end{center}%
        \end{minipage}
      }%
    \end{center}%
  \end{figure}%
}

%
% Comando para incluir imágenes en formato apaisado (sin marco).
\newcommand{\figuraApaisadaSinMarco}[5]{%
  \begin{figure}%
    \begin{center}%
    \includegraphics[angle=90,height=#1\textheight,#5]{#2}%
    \caption{#3}%
    \label{#4}%
    \end{center}%
  \end{figure}%
}
% Para las tablas
\newcommand{\otoprule}{\midrule [\heavyrulewidth]}
%
% Nuevo comando para tablas pequeñas (menos de una página).
\newcommand{\tablaSmall}[5]{%
 \begin{table}
  \begin{center}
   \rowcolors {2}{gray!35}{}
   \begin{tabular}{#2}
    \toprule
    #4
    \otoprule
    #5
    \bottomrule
   \end{tabular}
   \caption{#1}
   \label{tabla:#3}
  \end{center}
 \end{table}
}

%
% Nuevo comando para tablas pequeñas (menos de una página).
\newcommand{\tablaSmallSinColores}[5]{%
 \begin{table}[H]
  \begin{center}
   \begin{tabular}{#2}
    \toprule
    #4
    \otoprule
    #5
    \bottomrule
   \end{tabular}
   \caption{#1}
   \label{tabla:#3}
  \end{center}
 \end{table}
}

\newcommand{\tablaApaisadaSmall}[5]{%
\begin{landscape}
  \begin{table}
   \begin{center}
    \rowcolors {2}{gray!35}{}
    \begin{tabular}{#2}
     \toprule
     #4
     \otoprule
     #5
     \bottomrule
    \end{tabular}
    \caption{#1}
    \label{tabla:#3}
   \end{center}
  \end{table}
\end{landscape}
}

%
% Nuevo comando para tablas grandes con cabecera y filas alternas coloreadas en gris.
\newcommand{\tabla}[6]{%
  \begin{center}
    \tablefirsthead{
      \toprule
      #5
      \otoprule
    }
    \tablehead{
      \multicolumn{#3}{l}{\small\sl continúa desde la página anterior}\\
      \toprule
      #5
      \otoprule
    }
    \tabletail{
      \hline
      \multicolumn{#3}{r}{\small\sl continúa en la página siguiente}\\
    }
    \tablelasttail{
      \hline
    }
    \bottomcaption{#1}
    \rowcolors {2}{gray!35}{}
    \begin{xtabular}{#2}
      #6
      \bottomrule
    \end{xtabular}
    \label{tabla:#4}
  \end{center}
}

%
% Nuevo comando para tablas grandes con cabecera.
\newcommand{\tablaSinColores}[6]{%
  \begin{center}
    \tablefirsthead{
      \toprule
      #5
      \otoprule
    }
    \tablehead{
      \multicolumn{#3}{l}{\small\sl continúa desde la página anterior}\\
      \toprule
      #5
      \otoprule
    }
    \tabletail{
      \hline
      \multicolumn{#3}{r}{\small\sl continúa en la página siguiente}\\
    }
    \tablelasttail{
      \hline
    }
    \bottomcaption{#1}
    \begin{xtabular}{#2}
      #6
      \bottomrule
    \end{xtabular}
    \label{tabla:#4}
  \end{center}
}

%
% Nuevo comando para tablas grandes sin cabecera.
\newcommand{\tablaSinCabecera}[5]{%
  \begin{center}
    \tablefirsthead{
      \toprule
    }
    \tablehead{
      \multicolumn{#3}{l}{\small\sl continúa desde la página anterior}\\
      \hline
    }
    \tabletail{
      \hline
      \multicolumn{#3}{r}{\small\sl continúa en la página siguiente}\\
    }
    \tablelasttail{
      \hline
    }
    \bottomcaption{#1}
  \begin{xtabular}{#2}
    #5
   \bottomrule
  \end{xtabular}
  \label{tabla:#4}
  \end{center}
}



\definecolor{cgoLight}{HTML}{EEEEEE}
\definecolor{cgoExtralight}{HTML}{FFFFFF}

%
% Nuevo comando para tablas grandes sin cabecera.
\newcommand{\tablaSinCabeceraConBandas}[5]{%
  \begin{center}
    \tablefirsthead{
      \toprule
    }
    \tablehead{
      \multicolumn{#3}{l}{\small\sl continúa desde la página anterior}\\
      \hline
    }
    \tabletail{
      \hline
      \multicolumn{#3}{r}{\small\sl continúa en la página siguiente}\\
    }
    \tablelasttail{
      \hline
    }
    \bottomcaption{#1}
    \rowcolors[]{1}{cgoExtralight}{cgoLight}

  \begin{xtabular}{#2}
    #5
   \bottomrule
  \end{xtabular}
  \label{tabla:#4}
  \end{center}
}


















\graphicspath{ {./img/} }

% Capítulos
\chapterstyle{bianchi}
\newcommand{\capitulo}[2]{
	\setcounter{chapter}{#1}
	\setcounter{section}{0}
	\chapter*{#2}
	\addcontentsline{toc}{chapter}{#2}
	\markboth{#2}{#2}
}

% Apéndices
\renewcommand{\appendixname}{Apéndice}
\renewcommand*\cftappendixname{\appendixname}

\newcommand{\apendice}[1]{
	%\renewcommand{\thechapter}{A}
	\chapter{#1}
}

\renewcommand*\cftappendixname{\appendixname\ }

% Formato de portada
\makeatletter
\usepackage{xcolor}
\newcommand{\tutor}[1]{\def\@tutor{#1}}
\newcommand{\course}[1]{\def\@course{#1}}
\definecolor{cpardoBox}{HTML}{E6E6FF}
\def\maketitle{
  \null
  \thispagestyle{empty}
  % Cabecera ----------------
\noindent\includegraphics[width=\textwidth]{cabecera}\vspace{1cm}%
  \vfill
  % Título proyecto y escudo informática ----------------
  \colorbox{cpardoBox}{%
    \begin{minipage}{.8\textwidth}
      \vspace{.5cm}\Large
      \begin{center}
      \textbf{TFG del Grado en Ingeniería Informática}\vspace{.6cm}\\
      \textbf{\LARGE\@title{}}
      \end{center}
      \vspace{.2cm}
    \end{minipage}

  }%
  \hfill\begin{minipage}{.20\textwidth}
    \includegraphics[width=\textwidth]{escudoInfor}
  \end{minipage}
  \vfill
  % Datos de alumno, curso y tutores ------------------
  \begin{center}%
  {%
    \noindent\LARGE
    Presentado por \@author{}\\ 
    en Universidad de Burgos --- \@date{}\\
    Tutor: \@tutor{}\\
  }%
  \end{center}%
  \null
  \cleardoublepage
  }
\makeatother

\newcommand{\nombre}{Álvaro Delgado Pascual} %%% cambio de comando

% Datos de portada
\title{UBUAssistant\\Interacción por voz con la plataforma Moodle}
\author{\nombre}
\tutor{D. Raúl Marticorena Sanchez}
\date{\today}

\begin{document}

\maketitle


\newpage\null\thispagestyle{empty}\newpage


%%%%%%%%%%%%%%%%%%%%%%%%%%%%%%%%%%%%%%%%%%%%%%%%%%%%%%%%%%%%%%%%%%%%%%%%%%%%%%%%%%%%%%%%
\thispagestyle{empty}


\noindent\includegraphics[width=\textwidth]{cabecera}\vspace{1cm}

\noindent D. Raúl Marticorena Sánchez, profesor del departamento de Ingeniería Informática, área de Lenguajes y Sistemas Informáticos.

\noindent Expone:

\noindent Que el alumno D. \nombre, con DNI 71363793Z, ha realizado el Trabajo final de Grado en Ingeniería Informática titulado UBUAssistant interacción por voz con la plataforma Moodle. 

\noindent Y que dicho trabajo ha sido realizado por el alumno bajo la dirección del que suscribe, en virtud de lo cual se autoriza su presentación y defensa.


\begin{center} %\large
En Burgos, {\large \today}
\end{center}

\vfill\vfill\vfill

% para casos con solo un tutor comentar lo anterior
% y descomentar lo siguiente
Vº. Bº. del Tutor:\\[2cm]
D. Raúl Marticorena Sánchez


\newpage\null\thispagestyle{empty}\newpage




\frontmatter

% Abstract en castellano
\renewcommand*\abstractname{Resumen}
\begin{abstract}
	Los asistentes de voz son una nueva tecnología cuyo uso se ha disparado durante los útlimos años. Estos asistentes nos permiten realizar distintas tareas de una forma muy cómoda, interactiva y empleando poco tiempo. Tareas como consultar el tiempo, una búsqueda en Internet, la compra... en definitiva una gran variedad de tareas, en muchos campos distintos.
	
	En el ámbito de la educación, la enseñanza \textit{online} y las aulas virtuales también han tenido un importante desarrollo tecnológico, siendo Moodle un referente en este campo. Sin embargo, los asistentes de voz no se han utilizado en este área, a pesar de sus características y posibilidades que ofrecen para los estudiantes. Es por ello que ha surgido la idea de este proyecto, para dar a los usuarios de estas plataformas de educación una nueva forma de interactuar con ellas, de una forma más rápida y cómoda.
	
	UBUAssistant es una aplicación cliente desarrollada en Python utilizando el asistente de voz Mycroft. Esta aplicación permite que el usuario, mediante comandos de voz o a través de texto, obtenga información sobre el curso académico sin necesidad de conocer como funciona Moodle ni saber cómo tiene que navegar por la plataforma para obtener esta información.
\end{abstract}

\renewcommand*\abstractname{Descriptores}
\begin{abstract}
	Asistente de voz, Moodle, REST, intent, utterance, TTS, STT, skill
\end{abstract}

\clearpage

% Abstract en inglés
\renewcommand*\abstractname{Abstract}
\begin{abstract}
	Voice assistants are a new technology whose use has exploded over the last few years. These assistants allow us to perform varying tasks in an interactive and very comfortable way. Tasks such as check the weather, an Internet search, go shopping... in short a wide variety of tasks in many of fields.
	
	In the educational enviroment, online teaching and virtual classrooms have had an important technological development too, Moodle being a reference in this field. However, voice assistants have not been used in this area, despite of its features and the posibilities they bring to students. This is the reason that the idea of the project was conceived, to bring the users of these educational platforms a new way of interacting with them, in a quicker and easier way.
	
	UBUAssistant is a client application developed in Python using Mycroft as the voice assistant. This application allows the user to obtain information of the academic course through voice commands or text input, without the need of knowing what is Moodle or what steps he/she needs to do in order to retrieve this information.
\end{abstract}

\renewcommand*\abstractname{Keywords}
\begin{abstract}
	Voice assistant, Moodle, REST, intent, utterance, TTS, STT, skill
\end{abstract}

\clearpage

% Indices
\tableofcontents

\clearpage

\listoffigures

\clearpage

\listoftables
\clearpage

\mainmatter
\capitulo{1}{Introducción}

Este proyecto se basa en la utilización un asistente virtual para realizar diferentes tareas en una plataforma de Moodle, como es  UBUVirtual. La utilización de estas nuevas tecnologías permiten obtener información de la plataforma de Moodle de una forma más interactiva y cómoda, sin tener que navegar por menús y menús.

En esta memoria se van a ir explicando de forma ordenada los conceptos necesarios para entender el proyecto, las técnicas y herramientas utilizadas para su realización y los problemas y decisiones que se han surgido durante el desarrollo del mismo.
\capitulo{2}{Objetivos del proyecto}

Este apartado explica de forma precisa y concisa cuales son los objetivos que se persiguen con la realización del proyecto. Se puede distinguir entre los objetivos marcados por los requisitos del software a construir y los objetivos de carácter técnico que plantea a la hora de llevar a la práctica el proyecto.

\capitulo{3}{Conceptos teóricos}

A continuación se va a detallar qué es un asistente de voz, una \textit{Skill}, qué elementos tiene y cómo funcionan. También se explicará el funcionamiento de \textit{Mycroft}, que es el asistente de voz que se ha usado para el desarrollo del proyecto.

\section{Asistente de voz}

Un asistente de voz es un programa que actúa para un usuario, al que ayuda automatizando y realizando tareas con una interacción hombre-máquina mínima. La interacción entre el asistente y la persona debe ser una interacción lo más natural posible, es decir, la persona se comunica con la voz como es habitual y el asistente procesa, interpreta y responde de la misma forma.

\section{Skill}

Una \textit{Skill} es una aplicación diseñada para un asistente de voz. Estas \textit{skills} tienen diferentes elementos que hacen que funcionen, como la \textit{wake word}, las \textit{utterances}, los \textit{intents} y los \textit{dialogs}.

\subsection{Wake Word}

La \textit{wake word} es una palabra o palabras que tiene que decir el usuario para que el asistente comience a escuchar. En el asistente de voz que se ha usado para el proyecto, la \textit{wake word} es configurable, pero por defecto es "Hey Mycroft".

\subsection{Utterance}

Una \textit{utterance} es la frase completa que dice el usuario a continuación de la \textit{wake word} y que sirve como inicio del proceso de búsqueda de la \textit{Skill} correspondiente a la frase dicha. Un ejemplo de \textit{utterance} sería: "¿qué temperatura hace en Burgos?"

\subsection{Intent}

Un \textit{intent} es la parte de la frase o \textit{utterance} que el asistente de voz detecta como la parte necesaria para la invocación de una \textit{Skill}. Cada \textit{Skill} tiene un \textit{intent} asociado. Siguiendo con el ejemplo anterior, el intent sería: temperatura Burgos

\subsection{Dialog}

Este término es diferente al de otros asistentes. Un \textit{dialog} es una frase que dice el propio asistente, normalmente en respuesta a una pregunta por parte del usuario. Los \textit{dialog} son independientes de las \textit{Skill}, pero es normal que cada \textit{Skill} tenga un \textit{dialog}. Es el último paso cuando usamos una \textit{Skill}, pero no tiene por qué estar presente. Finalizando con el ejemplo, el \textit{dialog} sería: "Hay 11ºC en Burgos".
\capitulo{4}{Técnicas y herramientas}

\section{Metodología de desarrollo}
Para el desarrollo del proyecto se ha utilizado metodología ágil, en concreto SCRUM. Se ha adoptado una estrategia de desarrollo iterativa, con sprints bimensuales en los que se han realizado las entregas parciales. Las reuniones de estado del proyecto se han realizado semanalmente, a diferencia de lo normal en SCRUM que se realizan diariamente.
\section{Moodle}
\href{https://moodle.org/}{Moodle} es una plataforma de aprendizaje que permite a docentes y alumnos impartir y recibir clases a distancia, así como facilitar la gestión de los cursos. Es una herramienta de código abierto y por lo tanto cualquiera puede utilizarlo. En su página se pueden encontrar versiones de prueba como \href{https://school.moodledemo.net/}{Mount Orange School}
\section{PyQt5}
\href{https://doc.qt.io/qtforpython/}{PyQt5} es un binding (una adaptación de una biblioteca para que sea usada en otro lenguaje). Qt es un framework multiplataforma orientado a objetos para desarrollar programas que utilizan interfaces gráficas de usuario. Es software libre y de código abierto.
\section{requests}
\href{https://requests.readthedocs.io/es/latest/}{Requests} es una biblioteca para Python que permite enviar peticiones HTTP de forma muy sencilla.
\section{MiKTeX}
\section{TeXstudio}
\capitulo{5}{Aspectos relevantes del desarrollo del proyecto}

Este apartado pretende comentar los aspectos importantes y problemas que han surgido con la realización del proyecto, así como las decisiones de añadir o no más funcionalidad a la aplciación.

\section{Utilización de AmazonWebServices y Alexa para la realización del proyecto}

La idea inicial era utilizar los servicios web de Amazon (Amazon Lambda) y la consola de desarrollador (Amazon Developer) para crear la Skill de Alexa y hostearla en sus servidores.

Ya que nunca había usado nunca estos servicios ni había desarrollado ninguna Skill o aplicación asistente de voz similar, leyendo la documentación de Alexa encontré una guía de cómo crear tu primera skill.

\subsection{Amazon Web Services}

Los Amazon Web Services, AWS a partir de ahora, son un conjunto de servicios de computación en la nube. En concreto, de todos estos servicios se utilizó Amazon Lambda, que es una plataforma sin servidor basada en eventos, que es la base de los AWS. Lambda ejecuta código como respuesta a eventos y gestiona los recursos necesitados por el código ejecutado.

\subsection{Amazon Alexa}

Es un servicio de voz en la nube para dispositivos de Amazon y dispositivos de terceros que usan Alexa.

\section{Problemas con Alexa}


\section{Mycroft como alternativa}

Esta herramienta es de código abierto, por lo que tengo la posibilidad de ver como están implementados los diferentes servicios del asistente de voz, modificarlos, etc. Además, para crear una Skill de Alexa, estás restringido en cierto modo de qué puedes y qué no puedes hacer, problema que no existe con Mycroft que te da libertad total para crear lo que quieras. Tampoco tiene las restricciones de tamaño que me impedían continuar con el proyecto en Alexa, así que por todas estas características decidí volver a empezar pero usando Mycroft esta vez.

Pero con Mycroft no son todo ventajas. De momento esta aplicación no se puede instalar en Windows, teniendo que usar una máquina virtual de Ubuntu o distribución de Linux para usarla. Al ser una aplicación cliente el usuario necesita instalarla para usar la Skill.

\section{Aplicación cliente}

En el anterior apartado he comentado la libertad que te ofrece Mycroft. Pues bien, esta herramienta se ejecuta como una aplicación cliente, a diferencia de Alexa que necesitas usar los servicios de Amazon para usar la Skill. Y aunque Mycroft también depende de servicios de text-to-speech y speech-to-text, y por ello una conexión a internet, se puede crear una aplicación cliente usando estos servicios ya que la propia aplicación de Mycroft es cliente.

Así que debido a esto el proyecto tomó un camino distinto al inicial, teniendo la posibilidad de crear una aplicación gráfica más atractiva en general.

\section{Comunicación entre procesos}

En el inicio de esta nueva etapa surgió un problema que no había planteado cuando decidí usar Mycroft como alternativa a Alexa. Al ejecutar la aplicación gráfica para loguearse en Moodle e invocar una Skill, se crea un nuevo programa con esa Skill, por lo que la comunicación entre la aplicación gráfica y Skill se complicó. Como solución utilicé comunicación entre procesos mediante sockets.
\capitulo{6}{Trabajos relacionados}

El campo de este proyecto está muy inexplorado, aunque si que hay algún trabajo bastante similar a este proyecto.

\textbf{ALEXA SKILL VOICE INTERFACE FOR THE MOODLE LEARNING MANAGEMENT SYSTEM}\cite{meltonALEXASKILLVOICE} es un proyecto de una \textit{skill} para Alexa que permite a los usuarios acceder a anuncios de un curso, notas y eventos próximos.

Sin embargo, en el campo de los asistentes de voz existen múltiples proyectos en gran variedad de campos. En la \href{https://market.mycroft.ai/skills}{página} de Mycroft para conseguir nuevas \textit{skills} puedes por ejemplo \href{https://market.mycroft.ai/skills/3422768c-8f3a-45c4-ad47-a86b9eba6a02}{controlar las luces}, controlar \href{https://market.mycroft.ai/skills/b213f51b-c8aa-455e-80ec-f2489cf20eca}{las horas de buses} o \href{https://market.mycroft.ai/skills/5cff8320-4738-4e06-ac87-a5c3e6af5c25}{jugar a un juego}.

Además, el uso de los asistentes de voz se puede escalar a proyectos mucho mayores que los anteriores. Por ejemplo, la universidad de Saint Louis instaló 2300 \textit{Echo Dot}\cite{SLUAlexaProject}, que son los dispositivos para el asistente de voz Alexa, y crearon una \href{https://www.amazon.com/Saint-Louis-University-Ask/dp/B07YDNW9RN}{\textit{skill} personalizada} que ofrecía muchas posibilidades para los estudiantes, como preguntar por eventos y organizaciones, hora de cierre de la biblioteca o dónde se puede comer.

Las ventajas del proyecto de la universidad de Saint Louis respecto al mío son la facilidad de uso, ya que los estudiantes solo tienen que aprender a utilizar la \textit{skill} para tener acceso a la información que esta les provee, y que, siendo un proyecto mucho más grande hecho por la universidad, los usuarios tienen acceso a más funcionalidad, por ejemplo saber las cafeterías disponibles para comer. La desventaja principal sería el coste del proyecto, la instalación de los 2300 \textit{Echo Dot} rondaría los 80.500€ si cada dispositivo les cuesta 35€, que es lo que valen en la página de Amazon. Solo eso costaría 16 veces más que este proyecto.
\capitulo{7}{Conclusiones y Líneas de trabajo futuras}

\section{Conclusiones}

Una de las conclusiones que he sacado de este proyecto es que, aunque te encuentres con grandes problemas como fue el prototipo de skill para Alexa, en este tipo de proyectos es importante no abandonar y buscar otro camino, porque aunque parezca una pérdida de tiempo o algo problemático, puede ser que llegues a un mejor resultado que el esperado, como por ejemplo poder hacer una interfaz gráfica tras el cambio a Mycroft.

Otra conclusión a la que se ha llegado es que es importante dedicarle más tiempo a pensar qué es lo que quieres hacer que ponerse a hacer sin tener un plan, porque al final acabas dedicando más tiempo arreglando y cambiando cosas que el que dedicas cuando piensas las cosas.

Como conclusión final, se puede afirmar que se han cumplido los objetivos del proyecto. El producto final es una aplicación que permite al usuario interactuar con una plataforma de Moodle mediante voz y texto, configurar el asistente de voz y multitud más de funcionalidad.

Se han aplicado muchos de los conocimientos aprendidos durante el grado, como puede ser la interacción hombre-máquina, buenas prácticas de programación, defectos de diseño, refactorización de código, etc.

También se ha aprendido mucho en cuanto a la realización de interfaces gráficas, con las que no tenía apenas experiencia.

En definitiva, gracias a lo aprendido en la carrera se ha creado una aplicación funcional, útil y cumpliendo con los objetivos del proyecto. Ha sido una experiencia bastante agradable, aunque a veces te encuentres con problemas es muy satisfactorio solventarlos y te motiva para continuar con el proyecto y explorar un campo completamente desconocido para mí, como eran los asistentes de voz, es muy entretenido y se aprende mucho.

\section{Líneas de trabajo futuras}

Durante la realización del proyecto han surgido varias ideas sobre cómo se puede mejorar la aplicación.

Una de las mayores desventajas que tuvo cambiar de Alexa a Mycroft es que Mycroft no tiene soporte para Windows. Sin embargo, es posible ejecutar Mycroft desde Docker\cite{DockerSoftware2020}, que permite automatizar el despliegue de aplicaciones, por lo que se podría ejecutar la aplicación con Docker y utilizarla conectándose al contenedor desde Windows.

Otra posible continuación del desarrollo es guardar los datos que se obtienen de Moodle en ficheros locales de forma que se pueda utilizar la aplicación sin conexión a Internet. Usar la aplicación \textit{offline} reduciría bastante la funcionalidad de la aplicación, principalmente el servicio de STT, aunque gracias a que Mycroft es modular se pueden cambiar estos elementos. De hecho, el motor de TTS por defecto en Mycroft, Mimic\cite{Mimic} permite instalarse y utilizarse de forma local. En este \href{https://community.mycroft.ai/t/can-i-use-mycroft-offline/5306/5}{post del foro} de Mycroft se detalla bastante qué cambios son necesarios para utilizar Mycroft sin conexión.

Una de las formas más claras de mejorar la aplicación es aumentar la funcionalidad de la aplicación, permitiendo al usuario obtener más información de Moodle.

Por último se puede mejorar la interacción con el usuario. Esto implica mejorar la interfaz gráfica, que puede ser añadiendo funcionalidad o mejorar la claridad visual, y mejorar la interacción con Mycroft, de forma que las conversaciones sean más "humanas".


\printbibliography

\end{document}
